% GENERATED FILE. Edit canonical lessons, then run npm run generate:books.
% canonical-source-hash: fnv1a64:cab3b0fbfed8ba8a
% canonical-lessons: TE-C206-dorlu, TE-C206-jaru, TE-C206-veladu, TE-C206-tokku, TE-C206-kottu

\chapter{To roll, To slip, To hang, To tread on, To hit}
\label{ch:te-to-roll-to-slip-to-hang-to-tread-on-to-hit}
\hlchaptermodality{206}

\begin{chapteropening}
\textbf{By the end of this chapter:} \emph{I can say to roll, to slip, to hang, to tread on and to hit.}

Complete the last lesson of chapter 206: I can say to roll, to slip, to hang, to tread on and to hit.
\end{chapteropening}

\section[\texorpdfstring{\te{దొర్లు} (dorlu)}{దొర్లు (dorlu)}]{\te{దొర్లు} (dorlu) — to roll (over)}

\label{lesson:TE-C206-dorlu}



\begin{quote}
Before the new one: say the Telugu for to find out, then the Telugu for to borrow.
\end{quote}

\subsection*{You'll want to know: \te{దొర్లు}}
\textbf{\te{దొర్లు}} — \emph{dorlu} — \textquotedblleft{}to roll (over)\textquotedblright{}.

\begin{grammarlens}[title={What you've built}]
One more word for this chapter.
\end{grammarlens}

\subsection*{Your turn}
\begin{itemize}
\raggedright
  \item \emph{Say it:} \emph{dorlu}
  \item \emph{Say it:} \emph{dorlu}, once more
  \item \emph{Say it:} say \emph{dorlu}
  \item \emph{Recall:} say the Telugu for to find out, then the Telugu for to borrow, then say \emph{dorlu} again
\end{itemize}

\begin{tcolorbox}[breakable,colback=teal!4,colframe=teal!35!black,title={Before you move on}]
What does \te{దొర్లు} mean? (\textquotedblleft{}To roll (over)\textquotedblright{}.) Say it once more.
\end{tcolorbox}
\section[\texorpdfstring{\te{జారు} (jāru)}{జారు (jāru)}]{\te{జారు} (jāru) — to slip, to slide}

\label{lesson:TE-C206-jaru}



\begin{quote}
Before the new one: say the Telugu for to borrow, then the Telugu for to roll.
\end{quote}

\subsection*{You'll want to know: \te{జారు}}
\textbf{\te{జారు}} — \emph{jāru} — \textquotedblleft{}to slip, to slide\textquotedblright{}.

\begin{grammarlens}[title={What you've built}]
One more word for this chapter.
\end{grammarlens}

\subsection*{Your turn}
\begin{itemize}
\raggedright
  \item \emph{Say it:} \emph{jāru}
  \item \emph{Say it:} \emph{jāru}, once more
  \item \emph{Say it:} say \emph{jāru}
  \item \emph{Recall:} say the Telugu for to borrow, then the Telugu for to roll, then say \emph{jāru} again
\end{itemize}

\begin{tcolorbox}[breakable,colback=teal!4,colframe=teal!35!black,title={Before you move on}]
What does \te{జారు} mean? (\textquotedblleft{}To slip, to slide\textquotedblright{}.) Say it once more.
\end{tcolorbox}
\section[\texorpdfstring{\te{వేలాడు} (vēlāḍu)}{వేలాడు (vēlāḍu)}]{\te{వేలాడు} (vēlāḍu) — to hang, to dangle}

\label{lesson:TE-C206-veladu}



\begin{quote}
Before the new one: say the Telugu for to roll, then the Telugu for to slip.
\end{quote}

\subsection*{You'll want to know: \te{వేలాడు}}
\textbf{\te{వేలాడు}} — \emph{vēlāḍu} — \textquotedblleft{}to hang, to dangle\textquotedblright{}.

\begin{grammarlens}[title={What you've built}]
One more word for this chapter.
\end{grammarlens}

\subsection*{Your turn}
\begin{itemize}
\raggedright
  \item \emph{Say it:} \emph{vēlāḍu}
  \item \emph{Say it:} \emph{vēlāḍu}, once more
  \item \emph{Say it:} say \emph{vēlāḍu}
  \item \emph{Recall:} say the Telugu for to roll, then the Telugu for to slip, then say \emph{vēlāḍu} again
\end{itemize}

\begin{tcolorbox}[breakable,colback=teal!4,colframe=teal!35!black,title={Before you move on}]
What does \te{వేలాడు} mean? (\textquotedblleft{}To hang, to dangle\textquotedblright{}.) Say it once more.
\end{tcolorbox}
\section[\texorpdfstring{\te{తొక్కు} (tokku)}{తొక్కు (tokku)}]{\te{తొక్కు} (tokku) — to tread on, to pedal}

\label{lesson:TE-C206-tokku}



\begin{quote}
Before the new one: say the Telugu for to slip, then the Telugu for to hang.
\end{quote}

\subsection*{You'll want to know: \te{తొక్కు}}
\textbf{\te{తొక్కు}} — \emph{tokku} — \textquotedblleft{}to tread on, to pedal\textquotedblright{}.

\begin{grammarlens}[title={What you've built}]
One more word for this chapter.
\end{grammarlens}

\subsection*{Your turn}
\begin{itemize}
\raggedright
  \item \emph{Say it:} \emph{tokku}
  \item \emph{Say it:} \emph{tokku}, once more
  \item \emph{Say it:} say \emph{tokku}
  \item \emph{Recall:} say the Telugu for to slip, then the Telugu for to hang, then say \emph{tokku} again
\end{itemize}

\begin{tcolorbox}[breakable,colback=teal!4,colframe=teal!35!black,title={Before you move on}]
What does \te{తొక్కు} mean? (\textquotedblleft{}To tread on, to pedal\textquotedblright{}.) Say it once more.
\end{tcolorbox}
\section[\texorpdfstring{\te{కొట్టు} (koṭṭu)}{కొట్టు (koṭṭu)}]{\te{కొట్టు} (koṭṭu) — to hit, to beat}

\label{lesson:TE-C206-kottu}



\begin{quote}
Before the new one: say the Telugu for to hang, then the Telugu for to tread on.
\end{quote}

\subsection*{You'll want to know: \te{కొట్టు}}
\textbf{\te{కొట్టు}} — \emph{koṭṭu} — \textquotedblleft{}to hit, to beat\textquotedblright{}.

\begin{grammarlens}[title={What you've built}]
One more word for this chapter.
\end{grammarlens}

\subsection*{Your turn}
\begin{itemize}
\raggedright
  \item \emph{Say it:} \emph{koṭṭu}
  \item \emph{Say it:} \emph{koṭṭu}, once more
  \item \emph{Say it:} say \emph{koṭṭu}
  \item \emph{Recall:} say the Telugu for to hang, then the Telugu for to tread on, then say \emph{koṭṭu} again
\end{itemize}

\begin{tcolorbox}[breakable,colback=teal!4,colframe=teal!35!black,title={Before you move on}]
What does \te{కొట్టు} mean? (\textquotedblleft{}To hit, to beat\textquotedblright{}.) Say it once more.
\end{tcolorbox}
